\documentclass[../../main.tex]{subfiles}

\begin{document}

\chapter[AST, Validasi Kode, dan Penanganan Error]{Abstract Syntax Tree, Validasi Kode,\\ dan Penanganan Error}
\label{chap:ast-validasi}

\begin{outcomebox}
\textbf{Sub-CPMK:} \subcpmk{4.4}, \subcpmk{1.7}\\[3pt]
Setelah menyelesaikan bab ini, mahasiswa mampu:
\begin{enumerate}[label=\arabic*.,leftmargin=*]
  \item membuat AST untuk subset bahasa C;
  \item melakukan traversal AST dan validasi kode secara multi-pass;
  \item menguji keluaran pesan kesalahan sintaksis yang informatif;
  \item menguji dan men-\textit{debug} penanganan kesalahan semantik;
  \item menghasilkan pesan kesalahan yang menyebut baris, nama, dan
        jenis pelanggaran;
  \item menyelesaikan M5: validator dan proyek akhir.
\end{enumerate}
\end{outcomebox}

\section*{Materi Pokok}
\begin{itemize}[leftmargin=*]
  \item Desain dan konstruksi AST.
  \item Traversal AST dan visitor di C.
  \item Validasi kode: aturan dan multi-pass.
  \item Pengujian pesan kesalahan sintaksis dan semantik.
  \item Debug kesalahan dan diagnostik informatif.
\end{itemize}

\subfile{section-01}
\subfile{section-02}
\subfile{section-03}
\subfile{section-04}
\subfile{section-05}

\section*{Contoh}
\addcontentsline{toc}{section}{Contoh}

\begin{contoh}[Dump AST]
Untuk program
\begin{lstlisting}[language=C]
int main() {
    int a = 2 + 3;
    return 0;
}
\end{lstlisting}
dump AST kira-kira:
\begin{lstlisting}[language=bash]
PROGRAM
  FUNC main
    BLOCK
      DECL a
        BINARY +
          INT_LIT 2
          INT_LIT 3
      RETURN
        INT_LIT 0
\end{lstlisting}
Setiap simpul mencerminkan konstruksi program tanpa detail titik
koma atau kurung.
\end{contoh}

\begin{contoh}[Break di luar loop]
\begin{lstlisting}[language=C]
int main() {
    break;
    return 0;
}
\end{lstlisting}
Validator memasuki fungsi dengan \code{loop\_depth = 0}. Saat
mengunjungi \code{BREAK}, kedalaman masih nol, sehingga validator
melaporkan \code{break hanya valid di dalam loop}.
\end{contoh}

\begin{contoh}[Fungsi tanpa return]
\begin{lstlisting}[language=C]
int hitung(int x) {
    int y = x + 1;
}
\end{lstlisting}
Fungsi bertipe \code{int} tetapi tidak memiliki \code{return}.
Validator melaporkan \code{fungsi non-void harus mengembalikan
nilai} pada lokasi nama fungsi.
\end{contoh}

\begin{contoh}[Format diagnostik]
Keluaran
\begin{lstlisting}[language=bash]
tests/invalid/break_outside.sc:2:5: error: break hanya valid di dalam loop
\end{lstlisting}
memuat berkas, baris, kolom, tingkat, dan pesan. Editor dapat
melompat ke lokasi berdasarkan format ini.
\end{contoh}

\begin{contoh}[Menguji program tidak valid]
Setiap berkas di \file{tests/invalid} memiliki satu kesalahan utama.
Skrip \file{run-final-tests.bat} memastikan program valid diterima
dan program tidak valid ditolak. Pendekatan satu-kesalahan-per-berkas
memudahkan isolasi dan verifikasi.
\end{contoh}

\section*{Aktivitas Pembelajaran dan Praktikum}
\addcontentsline{toc}{section}{Aktivitas Pembelajaran dan Praktikum}

\begin{table}[H]
\centering
\caption{Rencana aktivitas Bab XVI}
\label{tab:ch16-aktivitas}
\begin{tabularx}{\textwidth}{@{}llX@{}}
\toprule
\textbf{Sesi} & \textbf{Bentuk} & \textbf{Kegiatan} \\
\midrule
1 & Ceramah & Desain AST dan traversal \\
2 & Live coding & Validator dan diagnostik \\
3 & Diskusi & Kualitas pesan kesalahan \\
4 & Praktikum & M5: proyek akhir dan demo \\
\bottomrule
\end{tabularx}
\end{table}

\subsection*{M5 / Proyek Akhir: Mini Compiler Front-End}
\textbf{Tujuan.} Menggabungkan lexer, parser, AST, type checker, dan
validator menjadi CLI yang bekerja.

\textbf{Keluaran.}
\begin{enumerate}[label=\arabic*.,leftmargin=*]
  \item CLI \code{scc} yang menerima satu berkas Subset-C.
  \item Pipeline lengkap: lexer, parser, type checker, validator.
  \item Minimal 40 kasus uji: valid dan tidak valid.
  \item Pesan kesalahan sintaksis dan semantik yang menyebut baris,
        nama, dan jenis pelanggaran.
  \item Log debug kesalahan yang ditemukan dan perbaikannya.
  \item Laporan proyek dan demo.
\end{enumerate}

\textbf{Langkah kerja.}
\begin{enumerate}[label=\arabic*.,leftmargin=*]
  \item Integrasikan seluruh modul.
  \item Tambahkan aturan validasi yang belum tercakup.
  \item Uji dengan program valid dan tidak valid.
  \item Ukur waktu dan penggunaan memori.
  \item Tulis laporan dan siapkan demo.
\end{enumerate}

\textbf{Rubrik.}
\begin{table}[H]
\centering
\caption{Rubrik Proyek Akhir}
\label{tab:ch16-rubrik}
\begin{tabularx}{\textwidth}{@{}lXr@{}}
\toprule
\textbf{Aspek} & \textbf{Indikator} & \textbf{Bobot} \\
\midrule
Fungsionalitas & Pipeline bekerja untuk program valid & 30\% \\
Integrasi & Modul terhubung dan konsisten & 20\% \\
Kualitas kode & Struktur, penamaan, komentar & 15\% \\
Pengujian & 40 kasus uji, kasus batas & 15\% \\
Dokumentasi & Laporan dan README & 10\% \\
Presentasi & Demo dan penjelasan & 10\% \\
\bottomrule
\end{tabularx}
\end{table}

\textbf{Definisi selesai.}
\begin{checklist}
  \item CLI memproses program valid tanpa kesalahan.
  \item CLI menolak program tidak valid dengan pesan jelas.
  \item Diagnostik menyebut baris, nama, dan jenis pelanggaran.
  \item 40 kasus uji tersedia.
  \item Log debug dan perbaikan terdokumentasi.
  \item Laporan dan demo selesai.
\end{checklist}

\section*{Latihan dan Refleksi}
\addcontentsline{toc}{section}{Latihan dan Refleksi}

\subsection*{Latihan}
\begin{enumerate}[label=\arabic*.,leftmargin=*]
  \item Jelaskan prinsip desain AST.
  \item Bandingkan simpul generik dan spesifik.
  \item Jelaskan pre-order dan post-order untuk AST.
  \item Rancang aturan validasi untuk \code{return}.
  \item Rancang aturan validasi untuk \code{main}.
  \item Jelaskan strategi pemulihan di setiap fase.
  \item Tuliskan format diagnostik yang baik.
  \item Mengapa jumlah pesan kesalahan perlu dibatasi?
\end{enumerate}

\subsection*{Refleksi}
\begin{reflectionbox}
\begin{enumerate}[label=\arabic*.,leftmargin=*]
  \item Bagian mana dari validator yang paling menantang?
  \item Apakah pesan kesalahan proyek Anda sudah informatif?
  \item Apa yang akan Anda perbaiki jika memulai proyek dari awal?
\end{enumerate}
\end{reflectionbox}

\section*{Asesmen}
\addcontentsline{toc}{section}{Asesmen}

Bab XVI memuat Project Final dengan bobot 20\% sesuai RPS. Asesmen
mencakup fungsionalitas, integrasi, kualitas kode, pengujian,
dokumentasi, dan presentasi.

\begin{table}[H]
\centering
\caption{Instrumen asesmen Bab XVI}
\label{tab:ch16-asesmen}
\begin{tabularx}{\textwidth}{@{}lXl@{}}
\toprule
\textbf{Instrumen} & \textbf{Indikator} & \textbf{Bobot} \\
\midrule
Project Final & Pipeline lengkap dan tes & 20\% \\
Laporan & Dokumentasi desain dan hasil & bagian proyek \\
Demo & Presentasi dan penjelasan & bagian proyek \\
\bottomrule
\end{tabularx}
\end{table}

Kriteria kelulusan: CLI memproses program valid, menolak program tidak
valid dengan diagnostik informatif, dan seluruh modul terintegrasi
dengan baik.

\section*{Checklist Pencapaian Kompetensi}
\addcontentsline{toc}{section}{Checklist Pencapaian Kompetensi}

\begin{checklist}
  \item Saya dapat merancang AST.
  \item Saya dapat melakukan traversal AST.
  \item Saya dapat menerapkan aturan validasi.
  \item Saya dapat mengklasifikasikan kesalahan.
  \item Saya dapat memulihkan dan melanjutkan analisis.
  \item Saya dapat menghasilkan diagnostik informatif.
  \item Saya dapat menguji program valid dan tidak valid.
  \item Saya sudah menyelesaikan M5/Project Final.
\end{checklist}

\section*{Rangkuman}
\addcontentsline{toc}{section}{Rangkuman}

Pada bab ini telah dibahas:
\begin{itemize}[leftmargin=*]
  \item desain dan konstruksi AST;
  \item traversal AST dan pola visitor di C;
  \item validasi kode dengan aturan dan multi-pass;
  \item klasifikasi, pemulihan, dan pelaporan kesalahan;
  \item pengujian pesan kesalahan sintaksis dan semantik, debug, dan
        proyek akhir.
\end{itemize}

Bab XVII akan menutup perjalanan dengan evaluasi, refleksi, dan
integrasi kompetensi. Kita akan memetakan kembali seluruh Sub-CPMK,
membahas blueprint ujian, dan merencanakan tindak lanjut.

\section*{Evaluasi}
\addcontentsline{toc}{section}{Evaluasi}

\begin{enumerate}[label=\arabic*.,leftmargin=*]
  \item Jelaskan prinsip desain AST.
  \item Bandingkan simpul generik dan spesifik.
  \item Jelaskan pre-order dan post-order.
  \item Rancang aturan validasi untuk \code{break}.
  \item Rancang aturan validasi untuk fungsi non-void.
  \item Jelaskan strategi pemulihan parser dan type checker.
  \item Tuliskan format diagnostik dan komponennya.
  \item Mengapa pengujian program tidak valid penting?
  \item Diberikan tiga program tidak valid, tunjukkan pesan kesalahan
        sintaksis/semantik yang menyebut baris, nama, dan jenis
        pelanggaran, lalu jelaskan perbaikannya.
\end{enumerate}

\subsection*{Kunci dan Rubrik Singkat}
\begin{itemize}[leftmargin=*]
  \item Soal 1--3: ketepatan konsep (70\%), contoh (30\%).
  \item Soal 4--6: rancangan aturan dan pemulihan (70\%),
        penjelasan (30\%).
  \item Soal 7--9: kualitas format, diagnostik, dan argumentasi
        (60\%), kelengkapan (40\%).
\end{itemize}


\ifSubfilesClassLoaded{%
  \bibliographystyle{apalike}%
  \bibliography{../../references}%
}{}

\end{document}
